% Transcription — fonds Alexandre Grothendieck, shelfmark 78, batch 5 (pages 81-100).
% Established from the facsimile archives/batches/78/batch-05.pdf (a mirror of
% https://grothendieck.umontpellier.fr/78.pdf), per the skill
% transcribe-grothendieck.
%
% Pass: Opus 5.5 (claude-opus-5-5), 2026-10-02 — first pass, unchecked against the pages by a human.
%
% Nineteen of the twenty pages are transcribed. Page 93 is an unrelated typed
% letter (CEREMADE, Paris, 10 November 1977) and is absent; its date is
% recorded in a note at page 94 because it bears on the folder's dating.
%
% The batch is one continuous run, in blue ink (pp. 81-91), then in a darker
% ink (pp. 92-95) and in black (pp. 96-100):
%   pp. 81-85 — affine coordinates on E attached to the base (s_0, s_1, s_{-1});
%     the matrices of sigma_0, sigma_1; the condition sigma_0^2 = id gives
%     alpha beta = (beta+2) beta = 0; Proposition (i)-(v) on beta = 0; the
%     "improper" case alpha = 0, beta = -2; Scholie: proper regular
%     realisations of Pi_infty are described by alpha in k.
%   pp. 86-89 — §3 fixed points of sigma_0, sigma_1, the centre; §4 eigenvalues
%     of u_V, alpha = zeta + zeta^{-1}, the cases alpha = ±2.
%   pp. 90-95 — §5 study by cases: I) alpha ≠ ±2, II) alpha = -2,
%     III) alpha = 2 (no centre); p. 92 is a page of scratch computations.
%   pp. 96-100 — §6 invariant quadratic forms f_alpha, circumscribed conic;
%     §7 discriminant of f_0; §8 discriminant of f; §9 alpha = 2: the conic is
%     a parabola on which u acts as v -> v+1.
% Equation numbers run (15)-(36); the numbering starts before the batch.
%
\documentclass[11pt,a4paper]{article}
\input{../preamble/grothendieck}

\folder{78}
\batch{5}
\pages{81}{100}
\dating{[à partir de 1977]}
\watermark{Édition de démonstration}
\foldertitle{[Polygones réguliers et polynômes cyclotomiques] : notes manuscrites (s.d.).}

\begin{document}

\page{81}
\note{la page s'ouvre au milieu d'un argument commencé avant ce lot : $E$, $k$, les sommets $s_i$, les automorphismes $\sigma_0$, $\sigma_1$ et les paramètres $\alpha$, $\beta$ sont introduits plus haut, et la numérotation des formules commence ici à (15).}

Nous utiliserons désormais la base affine $(s_0, s_1, s_{-1})$ pour identifier $E$ à $k^2$, donc

a) On choisit $s_0$ comme \emph{origine}

b) On choisit

(15)
\[
\begin{cases}
e_1 = s_1 - s_0 \\
e_2 = s_{-1} - s_0
\end{cases}
\]
comme vecteurs de base de l'espace des translations $V$ de $E$, et

c) on identifie $(x,y) \in k^2$ au point $s_0 + x e_1 + y e_2$, i.e. on écrit
\[
(x,y) = s_0 + x e_1 + y e_2 .
\]

Il y a lieu d'expliciter en termes de ces coordonnées toutes les données précédentes.
\[
\begin{cases}
s_0 = (0,0), \quad s_1 = (1,0), \quad s_2 = (0,1) \\
e_1 = (1,0) \quad e_2 = (0,1)
\end{cases}
\]
\note{la page porte bien « $s_2 = (0,1)$ » ; par (15) c'est $s_{-1}$ qui a ces coordonnées.}

Les \uncertain{automorphismes} affines $\sigma_0$, $\sigma_1$ sont donnés par leurs valeurs sur $s_0, s_1, s_{-1}$. On a

(\uncertain{16})
\[
\begin{cases}
\sigma_0 s_0 = s_1 \\
\sigma_0 s_1 = s_0 \\
\sigma_0 s_{-1} = s_2 = s_0 + (1+\alpha) e_1 + (1+\beta) e_2 \\
\sigma_0 e_1 = -e_1 \\
\sigma_0 e_2 = \alpha e_1 + (1+\beta) e_2
\end{cases}
\]
\note{dans la troisième ligne l'indice de $s$ est surchargé ; la lecture $s_{-1}$ est celle qu'impose $\sigma_0 s_{-1} = s_2$. Dans la dernière ligne, une première valeur de $\sigma_0 e_2$ est biffée et illisible, \struck{\ill{}}, et remplacée en dessous par $\alpha e_1 + (1+\beta) e_2$.}

\drawing{à droite, un croquis : la ligne brisée $s_{-1}, s_0, s_1, s_2$, avec les vecteurs $e_2$ de $s_0$ vers $s_{-1}$ et $e_1$ de $s_0$ vers $s_1$.}

[\emph{NB} on \uncertain{rappelle} :
\[
\begin{cases}
\beta = 0 \text{ ssi } s_2 - s_{-1} \parallel e_1 \\
\alpha = \zeta + \zeta^{-1} = 2\cos\theta
\end{cases}
\]
]

\page{82}
\emph{Ex.} Repr. régulière définie par $\zeta \in k$ ($\zeta \neq \pm 1$) : $s_i = (\zeta^i, \zeta^{-i})$, on trouve donc
\[
s_2 - s_{-1} = (\zeta^2, \zeta^{-2}) - (\zeta^{-1}, \zeta) = (\zeta^2 - \zeta^{-1}, \zeta^{-2} - \zeta)
\]
parallèle à
\[
s_1 - s_0 = (\zeta - 1, \zeta^{-1} - 1)
\]
de façon précise on a

(\ill{})
\[
\begin{cases}
(s_2 - s_{-1}) = (1 + \zeta + \zeta^{-1})(s_1 - s_0) = (1+\alpha)(s_1 - s_0) \\
\text{où} \quad \alpha = \zeta + \zeta^{-1}
\end{cases}
\]
\note{sous cette accolade, un départ biffé, \struck{\ill{} $\beta$}, griffonné.}

On a donc envie d'imposer la condition
\[
\beta = 0 .
\]

\emph{Proposition} Conditions équivalentes
\begin{enumerate}
  \item[(i)] $\beta = 0$
  \item[(ii)] $s_2 - s_{-1} = (1+\alpha)(s_1 - s_0)$
  \item[(ii')] $s_2 - s_{-1} \parallel s_1 - s_0$
  \item[(iii)] $\forall n \in \mathbb{Z}$ \quad $s_{n+3} - s_n = (1+\alpha)(s_{n+2} - s_{n+1})$
  \item[(iii')] $\forall n \in \mathbb{Z}$ \quad $s_{n+3} - s_n = (1+\alpha)(s_{n+2} - s_{n+1})$
  \item[(iv)] $\Leftrightarrow$ $\det \sigma_{0V} = -1$
  \item[(iv')] $\det u_V = 1$
  \item[(iv'')] $\operatorname{Tr} \sigma_{0V} = 0$
\end{enumerate}
\note{en (ii) et en (iii) un signe est biffé avant le « $=$ » (sans doute un $\parallel$ remplacé). Telle que la page l'écrit, (iii') répète (iii) mot pour mot ; d'après (ii') on attendrait $s_{n+3} - s_n \parallel s_{n+2} - s_{n+1}$.}

\struck{(\ill{} $\sigma_{0V}$ est une \emph{réflexion} i.e. est $\neq$ \ill{})}

\page{83}
i.e. on a

(17)
\[
\sigma_0 = \begin{pmatrix} -1 & \alpha & 1 \\ 0 & 1+\beta & 0 \\ 0 & 0 & 1 \end{pmatrix}
\]
\note{les deux coefficients de la deuxième colonne sont récrits sur des premières valeurs biffées, illisibles.}

(18)
\[
\sigma_0(x,y) = (1 - x + \alpha y, (1+\beta) y)
\]

(19)
\[
\begin{cases}
\sigma_1 s_0 = s_0 \\
\sigma_1 s_1 = s_{-1} \\
\sigma_1 s_{-1} = s_1 \\
\sigma_1 e_1 = e_2 \\
\sigma_1 e_2 = e_1
\end{cases}
\qquad \text{i.e.}
\]

(20)
\[
\sigma_1 = \begin{pmatrix} 0 & 1 & 0 \\ 1 & 0 & 0 \\ 0 & 0 & 1 \end{pmatrix}
\]

(21)
\[
\sigma_1(x,y) = (y,x)
\]

Il reste : écrire la condition $\sigma_0^2 = \mathrm{id}$. On trouve

(\ill{})
\[
\begin{cases}
\boxed{\alpha\beta = 0} \\
(1+\beta)^2 = 1 \quad \text{i.e.} \quad \boxed{\beta^2 + 2\beta = 0} \quad \text{i.e.} \quad \boxed{\beta(\beta+2) = 0}
\end{cases}
\]
\note{à droite de la première ligne, \struck{$(1+\beta)^2 = 1$} est biffé, puis récrit à la ligne suivante. Le numéro de la formule est surchargé.}

\uncertain{Donc} :

\emph{Proposition} Les représentations géom. affines \add{sur $k$} régulières de $\Pi_\infty$ sont \struck{définies} en corr. 1-1 corr. avec les couples $(\alpha,\beta)$ d'éléments de $k$ satisfaisant les conditions

(22)
\[
\alpha\beta = (\beta+2)\beta = 0
\]

\page{84}
(v) On n'est pas dans le « cas \struck{\ill{}} \add{impropre} » suivant \add{($\alpha = 0$, $\beta = -2$)}

(23)
\[
\begin{cases}
\text{car}\, k = p \neq 2, \quad \sigma_0 = \begin{pmatrix} -1 & 0 & 1 \\ 0 & -1 & 0 \\ 0 & 0 & 1 \end{pmatrix} \\
\sigma_0(x,y) = (1-x, -y)
\end{cases}
\]
\note{au début de la première ligne un mot est biffé, et après la matrice « i.e. » est suivi d'un \struck{$\sigma_0$} biffé ; au début de la seconde, \struck{\ill{}} est biffé.}
i.e. $\sigma_0$ = symétrie p.r. au milieu du segment $(s_0, s_1)$.

\emph{Dém.} L'équivalence de (i), (ii), (ii') résulte trivialt. de
\[
s_2 - s_{-1} = (1+\alpha) e_1 + \beta e_2
\]
et du fait que $e_1 = s_1 - s_0$ et $e_2$ forment une base de $V$. L'équivalence de (ii) et (iii), ou (ii') et (iii'), s'obtient en \struck{faisant} \ill{} que pour $n = -1$ on obtient les conditions de (ii) (ii'), et \uncertain{qu'inversement} (iii), (iii') s'obtiennent de (ii) (ii') en \uncertain{utilisant} $u^n$. Comme

(24)
\[
\begin{cases}
\det u_V = -\det \sigma_{0V} = 1 + \beta \\
\operatorname{Tr} \sigma_{0V} = \beta
\end{cases}
\]
\note{avant l'accolade, un début de formule est biffé, \struck{$\det \sigma_{0V}$ \ill{}}.}
ce qui prouve l'équivalence de (iv) (iv') (iv'') avec les autres conditions. Enfin, l'équivalence avec (v) est évidente par la \struck{propos} formule (22).

\emph{Remarques sur le cas impropre} À isomorphisme près, il y a un et un seul pour \ill{} $k$

\page{85}
donné. On a dans ce cas

(25)
\[
\begin{cases}
u(x,y) = (1-y, -x) \\
u^2(x,y) = (x+1, y-1)
\end{cases}
\]
i.e. $u^2$ est la translation par le vecteur $(1,-1)$. Donc en car. 0 $u^2$ est d'ordre infini, donc a fortiori $u$ l'est, donc la réalisation obtenue de $\Pi_\infty$ est d'ordre infini (on trouve qu'elle est fidèle…) Si car. $p > 0$, $u^2$ est d'ordre $p$, donc $u$ d'ordre $2p$, on trouve \struck{\ill{} $2p$} \add{donc} que la réalisation géom. impropre se factorise par $\Pi_{2p}$. On vérifie aisément que la réalisation géom. de $\Pi_{2p}$ est bel et bien fidèle, donc donne un $2p$-polygone régulier.

Avec la théorie projective des polygones réguliers, l'anomalie \ill{} des « polygones réguliers impropres » s'explique complètement, et on voit que du \uncertain{point de vue} projectif on n'a rien perdu en excluant le cas impropre, ce que nous ferons désormais.

\emph{Scholie} Le cas d'une réalisation géom. régulière \uncertain{propre} de $\Pi_\infty$ est décrit par $\alpha \in k$, qui peut être choisi arbitrairement…

\marginal{terminologie « invariant numérique »}

\page{86}
\section*{3. Points fixes de $\sigma_0$, $\sigma_1$, centre de la réal. géom.}

(26)
\[
\begin{cases}
\sigma_0(x,y) = (1 - x + \alpha y, y) \\
\sigma_1(x,y) = (y,x) \\
u(x,y) = \sigma_0 \sigma_1(x,y) = (1 - y + \alpha x, x)
\end{cases}
\]

\drawing{à droite, un croquis : la ligne brisée $s_{-1}, s_0, s_1, s_2$ et, en tirets, les axes de $\sigma_1$ (par $s_0$) et de $\sigma_0$ (entre $s_0$ et $s_1$).}

Points fixes de $\sigma_0$ :

(27)
\[
H_0 : \quad 2x - \alpha y - 1 = 0
\]
c'est une droite, sauf pour $p = 2$, $\alpha = 2$ (cas du carré en car. 2) où on trouve $H_0 = \emptyset$. Mais \add{passant} en coordonnées hom., on trouve encore une droite $z = 0$ : la droite à l'infini (Dans ce cas, $\sigma_0$ est la translation par \ill{} $e_1$)

Points fixes de $\sigma_1$

(28)
\[
H_1 : \quad x - y = 0
\]
c'est une droite

\emph{NB} Un automorphisme affine \add{d'un plan affine} \uncertain{tel que} l'ens. des points fixes \struck{est} une droite est appelé une \emph{réflexion affine}. On voit donc que les cas « propres » de la proposition précédente est caractérisé aussi par la condition

(iv''') $\sigma_0$ est une réflexion affine, ou \struck{\ill{}} une translation \struck{(dans le} \add{(ce dernier cas} dans le seul cas $p = 2$, $\alpha = 0$ i.e. carré en car. 2, et alors $\sigma_0$ est encore une réflexion projective…)

\page{87}
(iv$^{\text{iv}}$) $\sigma_{0V}$ est une réflexion, $\neq$ l'identité (\ill{})

\emph{NB} Dans le cas impropre, $\sigma_{0E}$ est la sym. p.r. à un pt (car. $\neq 2$), $\sigma_{0V}$ est la symétrie $v \mapsto -v$.

\emph{Points fixes sous $G$} = \emph{pts fixes sous $G^+$} = \emph{pts fixes sous $u$} :

(29)
\[
x_0 = y_0 = \frac{1}{2-\alpha}
\]
(qui existe ssi $\alpha \neq 2$, et est \ill{} unique)

\marginal{terminologie : le centre de la réalisation de $\Pi_\infty$ …}

\section*{4. Valeurs propres de $u_V$. Comparaison avec les \ill{}…}

(30)
\[
u_V = \begin{pmatrix} \alpha & -1 \\ 1 & 0 \end{pmatrix}
\]

(31)
\[
\begin{cases}
\det u_V = 1 \\
\operatorname{Tr} u_V = \alpha
\end{cases}
\]
\note{avant l'accolade, un \struck{$\det$} biffé.}

équation aux valeurs propres ($\zeta \in \bar{k}$)

(32)
\[
\zeta^2 - \alpha\zeta + 1 = 0
\]
discriminant

(33)
\[
\Delta(\alpha) = \alpha^2 - 4 = (\alpha - 2)(\alpha + 2)
\]
\ill{} d'un tel $\zeta$, on \ill{} trouve

(33)
\[
\alpha = \zeta + \zeta^{-1}
\]
\note{le numéro (33) est écrit deux fois.}

(\emph{NB} \ill{} $\zeta$, $\zeta'$ de cette équation satisfont

(34)
\[
\begin{cases}
\zeta + \zeta' = \alpha \\
\zeta\zeta' = 1
\end{cases}
\quad \text{i.e.} \quad \zeta' = \zeta^{-1} .
\]

\page{88}
On retrouve encore une fois que \struck{\ill{}} \ill{} ici coïncide avec ce que nous avions appelé (dans le cas des $n$-polygones \emph{réguliers} avec $n$ premier à la car. $p$) « l'invariant numérique » dudit. Cependant ici le cas d'une racine double

(35)
\[
\zeta = \zeta' \quad \text{i.e.} \quad \zeta^2 = 1 \quad \text{i.e.} \quad \zeta = \pm 1
\]
n'est pas exclu, il correspond au cas

(36)
\[
\alpha = \pm 2 \quad \text{ou encore} \quad \Delta(\alpha) = 0
\]
($\alpha = 2$ correspond à $\zeta = 1$, $\alpha = -2$ à $\zeta = -1$, les deux cas venant à se confondre ssi $p = 2$).

Cela tient compte précisément des 2 seules valeurs de $\zeta$ que nous avions été \emph{obligés} d'exclure, pour définir une repr. géom. \ill{} $\varphi^{\zeta} = (\varphi^{\zeta}_0, \varphi^{\zeta}_1)$ de $\Pi_\infty$ ! [Le cas $\zeta = 1$, $\alpha = 2$ est également \emph{distingué} comme le seul cas où il n'y ait pas de centre à distance finie….]
\note{les indices et exposants de $\varphi^{\zeta}$ sont mal formés ; la forme donnée est une lecture.}

\page{89}
\struck{5. Étude des cas \ill{}}

\emph{NB} On retrouve, pour $\alpha \neq \pm 2$ i.e. $\zeta \neq \pm 1$, \add{du moins si $\zeta \in k$}, la représentation $\varphi^{\zeta}$ de $\Pi_\infty$. (Donc les \emph{seuls} \emph{cas nouveaux} (sur $k$ alg. clos) sont précisément les cas $\alpha = \pm 2$ (i.e. $\zeta = \pm 1$) qu'on avait été amené précédemment à exclure. Nous allons les étudier séparément au § suivant. Notons dès à présent que, même si on décidait (à tort !) \ill{} \ill{} \ill{} ces deux cas, les considérant comme trop « exotiques », les \uncertain{dédaignant} en conséquence, ils y reviendraient indirectement, car a) ils sont valables pour tt corps, et \ill{} \add{\uncertain{même} sur les} anneaux, de base $k$, et \struck{b) ils} \uncertain{pas seulement} pour $k$ alg. clos, et b) ils permettent de suivre la variation de \struck{\ill{}} $\varphi(\alpha)$ en termes des paramètres $\alpha$, intrinsèques.

\page{90}
\section*{5) Étude par cas}

I) $\Delta(\alpha) \neq 0$ i.e. $\alpha \neq \pm 2$ (« cas ordinaire »)

Étendant au besoin le corps de base $k$ en $k(\zeta) = k(\zeta')$, et rapportant maintenant $E$ au système de coordonnées de centre $0$, ayant comme axes les droites propres de $u$ $D_\zeta$ et $D_{\zeta'}$, et pour lesquelles $s_0$ est de coordonnées $(1,1)$, on trouve par un argument déjà fait le cas-type $\Pi_\zeta$ (où $\zeta$ n'est pas nécessairement une racine de 1). Il y a une forme quadratique (\uncertain{non} \ill{}) invariante, canonique à un scalaire multiplicatif près, qu'on \uncertain{normalise} en exigeant que sa valeur en les sommets \add{de la représentation $\Pi_\zeta$} soit égale à 1… L'image \struck{$\Pi_\zeta$} \add{est un polygone} géométrique, dont le nombre de côtés, égal à l'ordre de $\zeta$, est (\struck{éventuellement} \add{quand il est} fini) premier à la caractéristique.

\struck{II) $\alpha = -2$}
\[
\begin{array}{l}
\text{\struck{$\sigma_0(x,y) = (1 - x - 2y, y)$}} \\
\text{\struck{$u_0(x,y) = (1 - y - 2x, x)$}} \\
\text{\struck{$u_0^2(x,y) = (1 - x - 2(1 - y - 2x),$}} \\
\text{\struck{$1 - y - 2x) = (-1 + 3x \ldots$}}
\end{array}
\]
\struck{On \ill{} ici $\alpha \neq 2$ (i.e. $p \neq 2$), de sorte que l'on a encore existence d'un centre $O$.}
\note{tout ce cas II est encadré et barré de traits obliques ; il est repris p.~94.}

\page{91}
II Considérons le cas $\alpha = \pm 2$, donc $\zeta = \pm 1$. On aura pour la partie homogène \add{$u_0$} de $u$
\[
u_0 = \varepsilon(\mathrm{id} + v_0), \qquad \varepsilon = \pm 1, \quad v_0^2 = 0
\]
\add{i.e.}
\[
v_0 = \varepsilon u_0 - \mathrm{id}
\]
De façon explicite
\[
\begin{array}{l}
u(x,y) = (1 - y + 2\varepsilon x, x) \\
u_0(x,y) = (-y + 2\varepsilon x, x) \\
v_0(x,y) = (x - \varepsilon y, \varepsilon(x - \varepsilon y)) = (x - \varepsilon y)(1, \varepsilon) \\
\phantom{v_0(x,y)} = a'(z)\, a
\end{array}
\]
\note{devant $a'(z)$ une lettre surchargée, illisible.}
i.e. \struck{\ill{}} $a' \in \check{V}$, $a \in V$,
\[
\begin{cases}
a'(x,y) = x - \varepsilon y \\
a = (1, \varepsilon)
\end{cases}
\]

\page{92}
\note{page de calculs de brouillon, d'une encre plus sombre ; on en donne les formules dans l'ordre de la page.}
\[
\begin{array}{l}
u^2 = (1 - x + \alpha(1 - y + \alpha x), 1 - y + \alpha x) \\
\phantom{u^2} = ((1+\alpha) + (\alpha^2 - 1)x - \alpha y, 1 + \alpha x - y)
\end{array}
\]
\struck{$= (1 + 3x + 2y, 1 \ldots$}
\[
3 + 3x - 3y, \; 1 + 2x - y
\]
\[
x = y \qquad 2x = \alpha y \qquad \alpha \neq 2 \qquad
\begin{cases}
\sigma_1(X,Y) = (Y,X) \\
\sigma_0(X,Y) = (-X + \alpha Y, Y) \\
u(X,Y) = (-Y + \alpha X, X)
\end{cases}
\]
$\alpha = -2$ \quad \struck{$u(X,Y) =$}
\[
u = -(1 + v), \qquad v = -1 - u
\]
\[
\begin{array}{l}
u(X,Y) = (-Y - 2X, X) \\
v(X,Y) = (X + Y, -(X + Y))
\end{array}
\qquad v^2 = 0
\]
\note{la ligne de $v$ est corrigée sur la page (des signes surchargés, une parenthèse biffée).}
\[
\begin{cases}
u^2 = 1 + 2v = \begin{pmatrix} 1 & \\ 0 & 1 \end{pmatrix} \ldots
\end{cases}
\]
\drawing{à gauche, deux axes et une droite oblique par l'origine.}
\[
\begin{cases}
\sigma_1(x,y) = (y,x) \\
\sigma_0(x,y)
\end{cases}
\qquad u_V = \varepsilon(1 + v), \quad \varepsilon = \pm 1
\]
\[
\begin{array}{l}
u(X,Y) = (-Y + 2\varepsilon X, X) = \varepsilon( \ldots \\
u_V = \varepsilon(\mathrm{id} + v) \qquad v = \varepsilon u_V - \mathrm{id}
\end{array}
\]
\[
v(X,Y) = (X - \varepsilon Y, \varepsilon(X - \varepsilon Y)) = (X - \varepsilon Y) \cdot (1, \varepsilon)
\qquad X = \varepsilon Y
\]
\note{$X - \varepsilon Y$ et $\varepsilon(X - \varepsilon Y)$ sont surmontés de $X'$ et $Y'$ ; suivent plusieurs essais biffés, dont \struck{$X' = \varepsilon Y'$} et \struck{$X' - \varepsilon Y' = 0$}.}
\[
\begin{array}{l}
(X + Y, -(X + Y)) \\
(X - Y, -X - Y)
\end{array}
\]

\page{94}
\note{le feuillet 93, non transcrit, est une lettre dactylographiée datée de Paris, 10 novembre 1977, adressée à A.~Grothendieck ; c'est un terme a quo pour ces pages, compatible avec la datation de l'inventaire.}

I) $\alpha \neq \pm 2$ (valeurs propres distinctes de $u_V$, existe centre)

\struck{Sur l'extension finie extension $k'$ de $k$} Si $k$ contient les racines $\zeta$, $\zeta' = \zeta^{-1}$ de l'équation aux valeurs propres, alors on est isom. à \struck{$\Pi_\zeta$} $\varphi_\zeta$ … Ordre de $u_V$ est égal à celui de $u$, ou \uncertain{encore} : celui de $\zeta$

\struck{III) $\alpha = -2$, $\alpha \neq 2$}

\struck{II) Cas $\alpha = \pm 2$ i.e. valeurs propres confondues}

\struck{II) $\alpha =$}

\struck{II) Cas $\alpha = -2$, $\alpha \neq 2$ il y a}
\struck{$u =$}
\note{ces quatre départs sont biffés, ainsi que la matrice qui suit $u =$, de lignes $(-2, -1, 1)$, $(1, 0, 0)$, $(0, 0, 1)$ ; son premier coefficient est surchargé.}

II) $\alpha = -2$ \qquad $\zeta = \zeta' = -1$ \qquad \struck{$u_V^2 = \mathrm{id}$} $u_V^2$ = transvection d'ordre $p$, $u_V$ d'ordre $2p$, $u_E$ d'ordre $2p$ si $p \neq 2$ (car $\exists$ centre)

\uncertain{Mais} contradiction si $p = 2$. En effet, on a alors $\alpha = -2 = 0$, or le cas $\alpha = 0$ est celui du \emph{carré} ($\zeta^2 = -1$ …)
\[
\begin{array}{l}
u(x,y) = (1 - y, x) \\
u^2(x,y) = (1 - x, 1 - y) \\
u^4(x,y) = (x, y) \quad \text{---}
\end{array}
\]
\note{devant $u(x,y)$, un \struck{$\sigma_0(x,y)$} biffé.}

Si $p \neq 2$, \add{i.e. $\alpha \neq 2$,} on trouve, en prenant
\[
\begin{cases}
X = x - \frac{1}{2-\alpha} \\
Y = y - \frac{1}{2-\alpha}
\end{cases}
\]
\note{le signe devant $\frac{1}{2-\alpha}$ est un pâté ; le moins est celui qu'impose le centre (29).}
\[
\begin{cases}
\sigma_0(X,Y) = (-X + \alpha Y, Y) \\
\sigma_1(X,Y) = (Y, X) \\
u(X,Y) = (-Y + \alpha X, X) = -(\mathrm{id} + a' \otimes a)(X,Y)
\end{cases}
\]
\[
a'(X,Y) = X + Y \qquad a = (1, -1)
\]

\page{95}
\drawing{à gauche, un croquis : les axes, la droite $D$ de direction $a$ portant $s_0$ et $-s_0$, une flèche $a$, et le point $-s_0' = -(s_0 + \frac{1}{2}a)$.}
\[
s_0 = \left(\frac{1}{2-\alpha}, \frac{1}{2-\alpha}\right) = \left(\frac{1}{4}, \frac{1}{4}\right)
\]
\note{devant la parenthèse, un \struck{$X$} biffé.}
\[
\begin{array}{l}
u = -(\mathrm{id} + v_0) \\
u^2 = \mathrm{id} + 2 v_0
\end{array}
\]
\[
\begin{cases}
u^{2n} = \mathrm{id} + 2n v_0 \\
u^{2n+1} = -(\mathrm{id} + (2n+1) v_0)
\end{cases}
\]
\[
s_{2n} = s_0 + 2n\, \underbrace{v_0(s_0)}_{(\frac12, -\frac12)} = s_0 + n \underbrace{(1,-1)}_{a} = s_0 + na
\]
\[
s_{2n+1} = -(s_0 + (2n+1) v_0(s_0)) = -s_0 - \left(n + \tfrac12\right) a
\]
\[
\phantom{s_{2n+1}} = \left(-s_0 - \tfrac12 a\right) - na
\]
\note{les deux signes moins de la dernière ligne sont des pâtés ; ils sont lus d'après le calcul.}

III) $\alpha = +2$ (Pas de centre) \qquad $\zeta = \zeta' = +1$

\struck{\ill{}} $u_V = \mathrm{id} + v_0$ \qquad $v_0 = a' \otimes a \neq 0$, $v_0^2 = 0$

$u_V$ d'ordre $p$

Je dis que $u_E$ d'ordre $p$, sauf si $p = 2$ (déjà traité, où on trouve $2p = 4$)

\struck{$u_{\hat E}$} Du \uncertain{point de vue} de l'enveloppe vectorielle
\[
\begin{array}{l}
u_E = \mathrm{id} + w \qquad w^3 = 0 \\
u_E^p = \mathrm{id} + w^p \qquad p \geq 3 \text{ --- } \ldots
\end{array}
\]
\note{dans la seconde ligne, un coefficient \struck{$p$} biffé devant $w^p$, puis quelques signes biffés illisibles ; après le tiret, un mot \ill{}. Au-dessus de la ligne suivante, \struck{$p = 2$ on trouve} est biffé.}

[Pour $p = 2$, on trouve $u_E^2 = \mathrm{id} + w^2 \neq \mathrm{id}$ car $w^2 \neq 0$]

\page{96}
\section*{6) formes quadratiques invariantes.}

On s'intéresse aux expressions
\[
f(x,y) = ax^2 + bxy + cy^2 + ux + vy + d
\]
polynomiales de degré $\leq 2$ sur $E$, qui correspondent aux formes quadratiques sur l'enveloppe vectorielle $\hat{E}$ de $E$. On cherche celles qui sont invariantes par $G = \operatorname{Aut} \Pi$, i.e. par $\sigma_0$, $\sigma_1$. Il y a les constantes, donc on est ramené à voir les conditions pour que \struck{$f$ soit inv.} \add{l'expression sans} terme constant
\[
f(x,y) = ax^2 + bxy + cy^2 + ux + vy
\]
soit invariante. L'invariance par $\sigma_1$ signifie que $f(x,y)$ est symétrique en $x$, $y$
\[
a = c, \quad u = v \quad \text{i.e.}
\]
\[
f(x,y) = a(x^2 + y^2) + bxy + u(x + y)
\]
et prenant en explicit l'invariance par $\sigma_0$, on trouve
\[
f = a f_\alpha
\]
\[
\boxed{f_\alpha(x,y) = (x^2 - x) + (y^2 - y) - \alpha xy}
\]
\note{la lettre $f$ de cette page est tracée surchargée, comme une capitale ; l'indice de $f_\alpha$ est lu $\alpha$.}

\emph{NB} Les fonctions $f$ \add{(p.l. de degré $\leq 2$)} nulles sur $S$ invariantes par $G$ sont les multiples de $f_\alpha$ ci-dessus.

\page{97}
Les fonctions $1$, $f_\alpha$ forment une base \add{sur $k$} de l'ens. des fonctions pl. de degré $\leq 2$ inv. par $G$. Parmi ces fonctions, $f_\alpha$ est caractérisée par les \struck{fonctions} conditions

a) $f_\alpha$ nulle sur $S$

b) La « partie homogène » $f_0$ de $f$ satisfait
\[
f_0(s_i - s_{i-1}) = 1 \qquad \forall i
\]
[\emph{NB}
\[
f_0(X,Y) = X^2 + Y^2 - \alpha XY
\]
donc
\[
f_0(1,0) = 1 \; ]
\]

\marginal{[\ill{} conique circonscrite (canonique) $C(f_\alpha)$ définie par $f = 0$ (indépendante de la normalisation). On vérifie que si $\alpha \neq 0, 1$ (i.e. cas des polyg. réguliers à $n \geq 5$ côtés) c'est la \emph{seule} conique circonscrite]}
\note{cette note marginale est écrite en remontant le long du bord gauche ; l'ordre des lignes est restitué.}

7) \struck{Le discri} \emph{Discriminant de $f_0$}

Le discriminant de $f_0$ est
\[
\delta(f_0) = 4 - \alpha^2 = (2 - \alpha)(2 + \alpha)
\]
\note{avant $4 - \alpha^2$, une première valeur biffée, illisible.}
qui est nul ssi $\alpha = \pm 2$ : la quadrique définie par $f_0$ est non dég. ssi on a $\alpha \neq \pm 2$, i.e. ssi on est dans le cas « ordinaire ». On aura
\[
f_0(X,Y) = (Y - \zeta X)(Y - \zeta' X)
\]
où $\zeta$, $\zeta'$ sont les solutions de l'équation aux valeurs propres. Les « \struck{droites} \add{directions} isotropes »

\page{98}
(sur une extension $k'$ de $k$ contenant $\zeta$, $\zeta'$) sont les droites homogènes \struck{\ill{}}
\[
Y - \zeta X = 0 \qquad Y - \zeta' X = 0
\]
qui sont aussi les droites propres sous $u_0$
\[
(\quad u_0(X,Y) = (-Y + \alpha X, X) \quad )
\]
[ce qui correspond au fait que $u_0$ est une \struck{(\ill{})} \emph{rotation} pour $f_0$, si $f_0$ est non dég. i.e. $\alpha \neq \pm 2$]

Dans le cas $\alpha = \pm 2$, ces deux directions sont confondues ; \struck{\ill{}} pour $\alpha = -2$, $\zeta = \zeta' = -1$, on trouve la direction
\[
X + Y = 0
\]
qui intervenait déjà comme privilégiée dans l'étude du « cas II \struck{\ill{}} ». Pour $\alpha = 2$, $\zeta = \zeta' = 1$, c'est la direction $X = Y$, \struck{\ill{}} on va l'examiner plus en détail plus bas.

8) Discriminant \add{direct} de $f$. La matrice de $f$ est
\[
\begin{pmatrix} 2 & -\alpha & -1 \\ -\alpha & 2 & -1 \\ -1 & -1 & 0 \end{pmatrix}
\]
\note{le dernier coefficient est surchargé ; la lecture $0$ est celle du terme constant nul de $f_\alpha$, et elle donne le discriminant qui suit.}
son discriminant est égal à
\[
\delta = 2\delta' \qquad \delta' = -(\alpha + 2)
\]

\page{99}
donc $f$ est régulière ssi on a $\alpha \neq -2$. Dans ce cas, la \emph{conique circonscrite} $C(f_\alpha)$ est non sing.

On notera que c'est le cas en particulier pour $\alpha = +2$, sauf si $k$ est de car. 2. Mais comme $f_0$ est dégénérée, la conique $C(f)$ est une \emph{parabole}.

Dans le cas $\alpha = -2$, on aura
\[
f(x,y) = (x^2 - x) + (y^2 - y) + 2xy = (x+y)^2 - (x+y)
\]
\[
= (x+y)(x+y-1) = \lambda(x,y)\,[\lambda(x,y) - 1]
\]
i.e. $f$ est le produit de deux fonctions linéaires affines ayant \uncertain{même} partie homogène $x + y$. Ainsi $C(f)$ est la réunion des deux droites parallèles \add{(distinctes)}, d'équations
\[
x + y = 0, \quad x + y = 1,
\]
qui étaient déjà en évidence dans l'étude du cas II (où on avait été amené à utiliser les coordonnées
\[
X = x - \tfrac14, \quad Y = y - \tfrac14, \quad x + y = X + Y + \tfrac12,
\]
donc les droites précédentes ont comme équations resp.
\[
X + Y = -\tfrac12, \quad X + Y = +\tfrac12 \; ).
\]

\page{100}
9) \emph{Retour sur le cas} $\alpha = 2$, $\alpha \neq -2$ (i.e. $\alpha = 2$ en car. $\neq 2$).

On a
\[
f(x,y) = (x - y)^2 - (x + y) = v^2 - w
\]
où
\[
\begin{cases}
v = x - y \\
w = x + y
\end{cases}
\quad \text{i.e.} \quad
\begin{cases}
x = \frac12(v + w) \\
y = \frac12(-v + w)
\end{cases}
\]
\note{dans la première formule, un exposant sur le premier $x$ est biffé.}

\drawing{à gauche, un croquis : les axes $x$, $y$ et les axes obliques $v$, $w$ par l'origine.}

Le calcul donne pour $u$, défini en coordonnées \uncertain{anciennes} par
\[
u(x,y) = (1 - y + 2x, x)
\]
en coordonnées $(v,w)$
\[
u(v,w) = (v + 1, 1 + 2v + w)
\]
en particulier, pour des points sur la parabole $C(f)$ d'équation $v^2 - w = 0$ i.e. $w = v^2$
\[
u(v, v^2) = (v + 1, (v+1)^2) .
\]
\note{après la parenthèse ouvrante, un début biffé, illisible.}
Donc si on paramètre $C(f)$ par la coordonnée $v$, alors $u|C(f)$ est décrit par la translation $v \mapsto v+1$ !
\note{la page s'arrête ici ; la suite du § 9, s'il y en a une, est au-delà de ce lot.}

\end{document}
